\section{Methods for the Japanese translation logs}
\label{app:enja}

% Moved from Sections 4.1, 4.3, 5.5 and 6.6 of the main body on 24 October 2026, without changing the procedure. The
% main body keeps the candidates, the segment rule, the quantity that answers question 4 and the rule for "exceeds".

\paragraph{Corpus.}
ENJA15 is a public corpus of English-to-Japanese translation processes recorded with Translog-II, in which Japanese was
typed with an input method~\cite{carl2016lrec,zou2026translog}. We use only the from-scratch translation sessions and
take the final translation recorded in each log as its text. The sessions were acquired on 12 October 2026 (UTC) from
the repository of the database that distributes them, at a fixed revision. The database states two licences on two of
its pages; both are non-commercial with attribution, and we follow the more restrictive one, which also requires
derived works to be shared alike~\cite{crittlicense}. The corpus has 76 from-scratch sessions by 39 translators on six
source texts, with 12 or 13 translators per source text. Thirty-seven translators translated two source texts and two
translated one, so translators and source texts are crossed. Sessions are excluded if the log or translation is missing,
if the translation is empty, or if the log contains neither an input-method event nor an inserted letter.

\paragraph{Forms.}
For ENJA15 the forms are R0j (no processing); R3j, which keeps the event type and the time and replaces each key event
by one of the categories romaji letter, conversion, commit, commit containing punctuation, BackSpace, Enter, navigation
and other, and removes every string (the key name, the committed text and the keys typed for it); R3j-t($w$), with the
press times of R3j rounded ($w = 10$, $100$, $1{,}000$\,ms); and R4($w$). The logs have no release events, so hold
times do not arise. Before the main computation we checked on all sessions, which include updates of uncommitted
strings, conversions, cancelled conversions and commits, that no row of R3j keeps a string of its session.

\paragraph{Attackers.}
In a translation, segments are split at the Japanese comma and full stop (U+3001, U+3002, U+FF0C and U+FF0E) and at
line breaks, and the value of a segment
is the standard number of keystrokes needed to type its reading with an input method, from the readings of a
morphological analyser and a romaji conversion table. In a log, among the input-method and insertion events, a letter
press adds 1, a BackSpace press subtracts 1, and a commit that contains punctuation closes the segment. Only A0 and A2
are run on these logs. A1, A1$'$, A3 and A4 are not, because they rely on rebuilding the final draft or on models of
English. A0 compares the number of letter presses minus BackSpace presses with the sum of the keystroke counts of the
candidate's readings. A2 aligns the segments as on KLiCKe, with the same groups and $\lambda$. Neither reads a string,
so they give the same values on R0j, R3j and R3j-t; in R0j the committed strings remain in the log and matching is
trivial for an attacker that reads them, so R0j is only a reference.

\paragraph{Analysis.}
The fourth question of the study asks how far above chance and above A0 the same kind of matching is on the Japanese
translation logs. The candidates for a log are translations of the same source text by other translators: all of them
($N$ = all), or nine drawn without replacement, 20 times per log ($N = 10$); identical translations tie. Accuracy, $1/N$
and the accuracy of A0 are averaged over the logs of each source text and then over the source texts. Since
translators and source texts are crossed, intervals come from 2,000 bootstrap resamples of translators with all their
sessions, without stratification, and differences are taken within each resample. The quantity that answers the
question is the accuracy of A2 on R3j, full texts and $N$ = all, with its differences from $1/N$ and from A0. We say
that matching ``exceeds'' chance and A0 only if the lower bounds of the intervals of both differences are above zero.
R3j-t($w$) and R4($w$) are reported with the same rule; since neither A0 nor A2 uses times, R3j-t($w$) gives the
values of R3j. Values per source text are only described, and R0j is a reference. ENJA15 is much smaller than KLiCKe, and its intervals are expected to be wide.
